\chapter{Introduction}
\label{ch:intro}

\section{Motivation}

Many scientific quantities are defined on the surface of a sphere, $S^2$, such as surface temperature, geomagnetic fields, atmospheric pressure, gravitational potential, and wind direction. In reality, we do not observe the field continuously. Instead, observations are finite, scattered, and often spatially irregular. Weather stations, satellite measurements, and survey sites are some examples of observation sources. The main problem here is to reconstruct a continuous field from sparse observations while also determining how trustworthy it is at each location.

Spherical harmonics are a natural basis for reconstructing fields on a sphere due to their properties of orthogonality, completeness for square-integrable functions on the sphere, and their relationship with the spherical Laplace-Beltrami operator. These properties allow the model to control the smoothness or roughness of the reconstructed field with Sobolev-type regularisation. The predecessor dissertation by Thomas (2025) developed regularised spherical harmonic reconstruction for scalar fields and vector fields using vector spherical harmonics. The previous work, which was applied to real global temperature and wind data, showed good reconstruction accuracy. It also showed that the choice of spherical harmonic degree $L$ affects the balance between model complexity and the available data \citep{thomas2025}. Regularisation in this case helps prevent overfitting, while different regularisation seminorms can effectively control model smoothness. The existing model performs well for point prediction, but uncertainty is still missing.

For example, users need to know whether the prediction is likely to be within $0.5^\circ\mathrm{C}$ or whether the actual value could differ by $10^\circ\mathrm{C}$. This is especially important in difficult areas with sparse observations, such as oceans and polar regions, as well as regions with strong or local physical variability. In this dissertation, the high-altitude Andes become an important finding, suggesting that uncertainty is important for identifying where the reconstruction is reliable and where it is difficult.

This concern is not only related to our temperature dataset. A reconstructed field may be used for downstream purposes such as deciding where more monitoring is needed, feeding climate information into another model, supporting agricultural or environmental analysis, and communicating predictions to non-specialists. The quality of the reconstructed field may therefore affect decision-making in real-life applications. Therefore, users need to know how much they should trust a prediction at a particular location. A single global RMSE only provides information on how well the model performs on average; it does not tell us about the model's performance at a specific location. These two questions can have very different answers when, as is generally true of spatial data, prediction error is heteroscedastic across the domain. In other words, prediction difficulty can vary across the domain, leading to heteroscedastic prediction errors. Therefore, global accuracy metrics alone are insufficient for understanding local reliability.

This motivates spatially adaptive uncertainty quantification. Conformal prediction is selected because it is model-agnostic and does not require a specific probability distribution for the model errors. It can be easily applied to an existing predictor, while its coverage guarantee is based on exchangeability rather than assuming that the base model is correctly specified.

So, why not Gaussian processes or Kriging? They are also established approach to spatial uncertainty quantification that can provide a predictive distribution at each location. However, they require assumptions about the underlying stochastic process, covariance structure, and choice of kernel. If the covariance or kernel is misspecified, the uncertainty estimates may become poorly calibrated. Standard GP methods can also become computationally expensive as the number of observations increases. Standard GP inference typically scales approximately as $O(n^3)$. Instead, conformal prediction provides a different approach that does not require the base predictor to be probabilistic and can work well with the existing spherical harmonic model. It does not require redesigning the reconstruction model and provides a formal marginal coverage guarantee under exchangeability. Therefore, conformal prediction is preferred for this dissertation.

\section{Contributions}

This dissertation extends the spherical harmonic reconstruction framework of \citet{thomas2025} with a spatially adaptive uncertainty quantification layer. It is built by adapting the adaptive geodesic conformal prediction framework of \citet{amirishahbazi2026}, which uses a manifold-valued response in regression problems (e.g.\ a direction on $S^2$). However, the climate reconstruction problem is scalar, making this adaptation different. This adaptation is novel because the geodesic structure of the sphere is used to define a spatially local difficulty estimator $\hat\sigma(x)$ over the input domain. The estimator is obtained using cross-validated $k$-nearest-neighbour regression on residuals based on great-circle distance, and the conformal scores for the scalar output are normalised by this estimator. The contributions of this dissertation are separated into five parts:

\begin{enumerate}
    \item \textbf{Methodological adaptation} - Move geodesic information from the manifold response to the spherical input locations for scalar temperature (\cref{sec:adaptation}).
    \item \textbf{New uncertainty layer} - Add local $\hat{\sigma}(x)$ estimation and adaptive conformal calibration to the existing spherical harmonic model.
    \item \textbf{Empirical validation} - Show better coverage across regions with different levels of difficulty while maintaining similar interval widths.
    \item \textbf{Robustness} - Test different $L$ values using two hyperparameter-selection criteria for $\lambda$.
    \item \textbf{Geographic interpretation} - Show that high uncertainty is concentrated in sparse or physically difficult regions, especially the Andes.
\end{enumerate}

\section{Dissertation structure}

\Cref{ch:background} introduces the mathematical background for spherical harmonics, Sobolev-regularised least squares and summarises the predecessor work. It also introduces split conformal prediction and describes adaptive and geometry-aware conformal prediction. 

\Cref{ch:data} describes the dataset and pre-processing pipeline. It also describes how the observations and the data-splitting procedure are prepared.

\Cref{ch:methodology} presents the methodology, main algorithm and experimental design, with an explanation of the base spherical harmonic reconstruction model, local difficulty estimator $\hat{\sigma}(x)$, standard and adaptive conformal prediction procedures, and a description of how the original manifold-response conformal framework is adapted to the scalar-response temperature problem. 

\Cref{ch:results} presents results on hyperparameter selection, standard-versus-adaptive coverage and width comparisons, conditional coverage diagnostics, degree-robustness study, and geographic uncertainty analysis.

\Cref{ch:discussion} discusses the findings in the wider conformal prediction literature and the limitations of the approach. It covers the adaptive method in relation to previous conformal prediction work. It also discusses alternative uncertainty methods such as Gaussian Processes and Kriging, broader practical implications, and limitations of the current approach.

\Cref{ch:conclusion} provides an overall conclusion about the dissertation and possible future extensions. In particular, considering the application of the framework to vector-valued wind fields which could be applied more directly. 

